\documentclass[paper=a4, fontsize=11pt]{scrartcl}

\usepackage[T1]{fontenc}
\usepackage{mathptmx}
\usepackage[english]{babel}
\usepackage{amsmath}
\usepackage{booktabs}
\usepackage{geometry}
\usepackage{hyperref}
\usepackage{setspace}
\geometry{a4paper, margin=2.2cm}

\onehalfspacing
\setlength\parindent{0pt}

\newcommand{\horrule}[1]{\rule{\linewidth}{#1}}

\title{
\normalfont \normalsize
\textsc{Universidad Panamericana, Graphical Simulation} \\ [20pt]
\horrule{0.5pt} \\[0.35cm]
\huge Green Grid Predictive Digital Twin \\
\horrule{2pt} \\[0.45cm]
}

\author{Emiliano Hinojosa Guzm\'an \\ Diego Am\'in Hernandez Pallares \\ Jos\'e Salcedo Uribe}
\date{\normalsize May 26, 2026}

\begin{document}

\maketitle

\section{Project Goal and Pitch}

The original Green Grid simulator used an ideal solar curve, so every clear day
looked like a perfect sine wave. That was useful for testing energy-flow code,
but it did not answer the real business question: how would a neighborhood with
different household types, wealth levels, solar adoption, batteries, and
dispatch strategies behave under realistic weather?

Our final project replaces the synthetic solar curve with a predictive solar
model trained from real San Francisco 2006 production and weather data. The
model predicts solar capacity factor from weather and time inputs, then the
digital twin scales that prediction by each household's installed PV capacity.
The dashboard turns the result into a decision story: who imports energy, who
exports energy, when batteries help, and which segments benefit economically.

\section{Dataset, Cleaning, and Data Dictionary}

The ML pipeline is implemented in \texttt{simulation/ml/prepare\_data.py}. It
uses actual distributed PV production, day-ahead and four-hour-ahead forecasts,
and NSRDB weather observations. Actual production is resampled from 5-minute to
15-minute intervals, matching the simulator tick. Weather is shifted by
\(-8\) hours to align the NSRDB timestamps with local solar production, numeric
weather fields are time-interpolated, and categorical weather flags are
forward-filled. The target is:

\[
\texttt{capacity\_factor} =
\operatorname{clip}\left(\frac{\texttt{actual\_mw}}{33.0}, 0, 1\right)
\]

The prepared base dataset contains 35,007 rows from
\texttt{2006-01-01 00:00} through \texttt{2006-12-31 15:30}. Rows are split
chronologically: January--September for training, October--November for
validation, and December for testing. The actual-vs-GHI correlation is 0.9498,
and no rows have production when GHI is zero, which supports the timestamp
alignment decision.

\begin{table}[h]
\centering
\small
\begin{tabular}{llrrrr}
\toprule
\textbf{Column} & \textbf{Meaning} & \textbf{Min} & \textbf{Max} & \textbf{Mean} & \textbf{Std.} \\
\midrule
\texttt{actual\_mw} & measured solar MW & 0.00 & 29.40 & 5.10 & 7.60 \\
\texttt{capacity\_factor} & normalized target & 0.00 & 0.89 & 0.15 & 0.23 \\
\texttt{da\_mw} & day-ahead forecast MW & 0.00 & 28.60 & 6.10 & 8.82 \\
\texttt{ha4\_mw} & four-hour forecast MW & 0.00 & 26.60 & 5.09 & 7.40 \\
\texttt{temperature\_c} & air temperature & 1.40 & 32.10 & 13.95 & 4.62 \\
\texttt{relative\_humidity\_pct} & humidity percentage & 15.37 & 100.00 & 77.97 & 16.97 \\
\texttt{dhi} & diffuse irradiance & 0.00 & 517.00 & 60.01 & 89.56 \\
\texttt{dni} & direct irradiance & 0.00 & 993.00 & 236.55 & 351.57 \\
\texttt{ghi} & global irradiance & 0.00 & 1035.00 & 199.27 & 289.01 \\
\texttt{solar\_zenith\_angle} & sun angle & 14.53 & 165.58 & 89.64 & 37.21 \\
\texttt{wind\_speed} & wind speed & 0.30 & 9.50 & 2.59 & 1.17 \\
\texttt{pressure} & atmospheric pressure & 1001.00 & 1034.00 & 1018.52 & 4.94 \\
\texttt{cloud\_type} & NSRDB cloud category & 0.00 & 10.00 & 2.24 & 2.79 \\
\bottomrule
\end{tabular}
\caption{Core data dictionary for the base prepared ML dataset.}
\end{table}

\section{Feature Engineering and Model Selection}

The active model is a from-scratch linear regression trained by batch gradient
descent in \texttt{simulation/ml/linear\_regression.py}. We do not use
scikit-learn or any other ML library. Features are z-score normalized using
training rows only, then the same normalization is applied to validation and
test rows.

Categorical \texttt{cloud\_type} is one-hot encoded because cloud categories
are labels, not continuous magnitudes. Hour of day and day of year are encoded
with sine/cosine pairs so midnight and year boundaries do not create artificial
numeric jumps. The final deployed model uses \texttt{weather\_time}: weather,
irradiance, one-hot cloud type, and cyclic time features.

\begin{table}[h]
\centering
\begin{tabular}{lrrr}
\toprule
\textbf{Feature group} & \textbf{Validation \(R^2\)} & \textbf{Test \(R^2\)} & \textbf{Test RMSE} \\
\midrule
\texttt{time\_only} & 0.4941 & 0.0880 & 0.145811 \\
\texttt{irradiance\_only} & 0.8841 & 0.8552 & 0.058091 \\
\texttt{weather\_only} & 0.8895 & 0.8601 & 0.057100 \\
\texttt{weather\_time} & 0.8961 & 0.8750 & 0.053986 \\
\texttt{forecast\_only} & 0.8725 & 0.8416 & 0.060766 \\
\bottomrule
\end{tabular}
\caption{Linear-regression comparison using the same split and hyperparameters.}
\end{table}

\texttt{weather\_time} was selected because it achieved the best test
performance among the compared groups and can be rebuilt inside the simulator
from timestamp and weather data. \texttt{forecast\_only} is a useful benchmark,
but it would require a forecast lookup during simulation, so it is not the
deployed model. The saved active model reached test \(R^2 = 0.8750\) and test
RMSE \(= 0.053986\) capacity factor.

\section{Simulation and Dashboard Integration}

The trained coefficients are stored in
\texttt{simulation/ml/model\_coefficients.json}. During the simulation,
\texttt{components/weather.py} loads the same shifted 15-minute weather
features used in training. \texttt{components/panel.py} calls
\texttt{SolarModel.predict\_kw()}, which predicts a capacity factor, clamps it
to \([0, 1]\), and multiplies it by the household's PV size. A physics guard
returns zero generation when \texttt{ghi <= 1}, preventing night generation.

The neighborhood simulation contains roughly 100 households across studio,
small, and large home types; low, middle, high, and luxury wealth levels; and
different energy strategies. Each tick updates weather, solar generation,
appliance-based load, inverter behavior, battery state of charge, grid imports,
grid exports, and savings. Raw outputs are written to
\texttt{simulation/output/}; \texttt{data\_pipeline.py} then publishes
dashboard-ready CSVs to \texttt{dashboard/data/}.

The dashboard is the storytelling layer. It shows production vs. consumption,
the duck curve, household type and wealth comparisons, economics, adoption,
battery/grid behavior, and an ML Data tab. The ML tab provides the raw-data
visual evidence required by the assignment: actual production vs. GHI, daily
production across the year, base vs. alternate data differences, a sample week
of actual vs. forecasts, and cloud-type distribution.

\section{Insights, Challenges, and Improvements}

The feature comparison confirmed the original intuition: time alone is not
enough to model real solar behavior, while irradiance explains most of the
target. Adding weather, cloud type, and cyclic time gives the best deployed
model and keeps inference practical. In the simulation, this matters because
cloudy periods reduce production, which changes battery charge, grid imports,
exports, and household savings.

The main technical challenge was keeping training and inference consistent.
If the model is trained on shifted 15-minute weather but the simulator uses
unshifted or hourly weather, the digital twin quietly becomes invalid. The
final architecture fixes that by making \texttt{prepare\_data.py} the single
source for the weather-feature lookup and model feature expansion.

Future improvements should include multi-year weather data, household
occupancy patterns, transformer-level grid constraints, battery degradation,
and Monte Carlo runs over adoption and weather assumptions. These would turn
the current deterministic pitch into a stronger risk-analysis tool.

\section{Development Log and Reflections}

\textbf{Development log.} Emiliano focused on the simulation and ML integration:
data preparation, timestamp alignment, feature engineering, from-scratch linear
regression, and replacing the sine solar model with the trained predictor.
Jos\'e focused on dashboard structure and visual storytelling, especially the
tabs that compare production, consumption, economics, adoption, and ML data.
Diego focused on documentation, report framing, and keeping the final pitch
understandable for a mixed technical and business audience.

\textbf{Personal reflections.} Emiliano's main lesson was that data alignment
is part of the model, not a side task; the \(-8\) hour weather shift changed
the credibility of the whole project. Jos\'e's main lesson was that dashboard
charts need a clear decision purpose, not just visual variety. Diego's main
lesson was that the best technical result still needs a business narrative:
the model matters because it changes the neighborhood energy story.

\end{document}
